OB\documentclass[../../main.tex]{subfiles}
\begin{document}

{\bf [解答]} $S$を，$xy$平面上の
\begin{align}
  &|x|\le\dfrac{1}{2}&|y|\le\dfrac{1}{2}
\end{align}
とする．P$(a,b)$とすれば，$T(\mathrm{P})$は
\begin{align}
  &|x-a|\le\dfrac{1}{2}&|y-b|\le\dfrac{1}{2}
\end{align}
で表される．対称性から，
\begin{align}
  a,b\ge0\label{eq:1}
\end{align}
で考える．

\begin{figure}[htb]
  \centering
  \begin{tikzpicture}[x=2.5cm, y=2.5cm, >=stealth, every node/.style={font=\small}]
    \coordinate (O) at (0, 0);

    % 座標軸
    \draw[->, thick, gray!80!black] (-0.9, 0) -- (1.8, 0) node[right] {$x$};
    \draw[->, thick, gray!80!black] (0, -0.9) -- (0, 1.8) node[above] {$y$};
    \node[below left=1pt, font=\footnotesize] at (O) {$O$};

    % 正方形 S: [-0.5, 0.5] x [-0.5, 0.5]
    \draw[line width=1.2pt, gray!85!black] (-0.5, -0.5) rectangle (0.5, 0.5);
    \node[below right=1pt, font=\footnotesize] at (0.5, 0) {$\frac{1}{2}$};
    \node[below left=1pt, font=\footnotesize] at (-0.5, 0) {$-\frac{1}{2}$};
    \node[above left=1pt, font=\footnotesize] at (0, 0.5) {$\frac{1}{2}$};
    \node[below left=1pt, font=\footnotesize] at (0, -0.5) {$-\frac{1}{2}$};
    \node[above left=1pt, font=\footnotesize, gray!85!black] at (-0.5, 0.5) {$S$};

    % 点 P(a, b): a = 0.65, b = 0.65
    \coordinate (P) at (0.65, 0.65);
    % 正方形 T(P): [a-0.5, a+0.5] x [b-0.5, b+0.5] = [0.15, 1.15] x [0.15, 1.15]
    \draw[line width=1.1pt, dashed, blue!85!black] (0.15, 0.15) rectangle (1.15, 1.15);
    \node[above right=1pt, font=\footnotesize, blue!85!black] at (1.15, 1.15) {$T(P)$};

    % 共通部分: [0.15, 0.5] x [0.15, 0.5]
    \fill[pattern=north east lines, pattern color=blue!75!black] (0.15, 0.15) rectangle (0.5, 0.5);
    \draw[line width=1.3pt, blue!85!black] (0.15, 0.15) rectangle (0.5, 0.5);

    % 点 P
    \node[circle, fill=red!85!black, inner sep=1.4pt] at (P) {};
    \node[above right=1pt, font=\footnotesize, text=red!85!black] at (P) {$P(a, b)$};

    % 共通部分の境界値 a-1/2, b-1/2
    \draw[densely dotted, thick, gray!70!black] (0.15, 0) node[below=2pt, font=\footnotesize] {$a-\frac{1}{2}$} -- (0.15, 0.15) -- (0, 0.15) node[left=2pt, font=\footnotesize] {$b-\frac{1}{2}$};
  \end{tikzpicture}
  \caption{正方形 $S$ と $T(P)$ の共通部分}
  \label{fig:square_intersection}
\end{figure}

これらが共通部分を持つとき，
\begin{align}
  &
  \begin{cases}
    a-1/2\le1/2 \\
    b-1/2\le1/2
  \end{cases}\nonumber\\
  \Longleftrightarrow
  &
  \begin{cases}
    a\le 1 \\
    b\le 1
  \end{cases}\label{eq:2}
\end{align}
が条件で，このもとで長方形領域
\begin{align}
  \begin{cases}
    a-1/2\le x\le1/2 \\
    b-1/2\le y\le1/2
  \end{cases}
\end{align}
が共通部分である．この面積は
\begin{align}
  &\left(\dfrac{1}{2}-\left(a-\dfrac{1}{2}\right)\right)\left(\dfrac{1}{2}-\left(b-\dfrac{1}{2}\right)\right) \\
  =&(1-a)(1-b)
\end{align}
である．従って，求める条件は
\begin{align}
  \dfrac{1}{2}\le (1-a)(1-b)\label{eq:3}
\end{align}
である．以上\cref{eq:1}，\cref{eq:2}，\cref{eq:3}からPの存在範囲は下図斜線部(境界含む)．

\begin{figure}[htb]
  \centering
  \begin{tikzpicture}[x=3.0cm, y=3.0cm, >=stealth, every node/.style={font=\small}]
    \coordinate (O) at (0, 0);

    % 存在領域の斜線塗り (境界含む)
    \fill[pattern=north east lines, pattern color=blue!75!black, domain=0:0.5, samples=60]
    plot (\x, {1 - 1/(2*(1-\x))}) --
    plot (0.5-\x, {-1 + 1/(2*(1-(0.5-\x)))}) --
    cycle;
    \fill[pattern=north east lines, pattern color=blue!75!black, domain=-0.5:0, samples=60]
    plot (\x, {1 - 1/(2*(1+\x))}) --
    plot (-\x-0.5, {-1 + 1/(2*(1+(-\x-0.5)))}) --
    cycle;

    % 座標軸
    \draw[->, thick, gray!80!black] (-0.8, 0) -- (0.9, 0) node[right] {$a$};
    \draw[->, thick, gray!80!black] (0, -0.8) -- (0, 0.9) node[above] {$b$};
    \node[below left=1pt, font=\footnotesize] at (O) {$O$};

    % 境界線 (太線)
    \draw[line width=1.3pt, blue!85!black, domain=0:0.5, samples=60]
    plot (\x, {1 - 1/(2*(1-\x))});
    \draw[line width=1.3pt, blue!85!black, domain=0:0.5, samples=60]
    plot (\x, {-1 + 1/(2*(1-\x))});
    \draw[line width=1.3pt, blue!85!black, domain=-0.5:0, samples=60]
    plot (\x, {1 - 1/(2*(1+\x))});
    \draw[line width=1.3pt, blue!85!black, domain=-0.5:0, samples=60]
    plot (\x, {-1 + 1/(2*(1+\x))});

    % 軸上の交点 (+-1/2, 0), (0, +-1/2)
    \node[circle, fill=blue!85!black, inner sep=1.2pt] at (0.5, 0) {};
    \node[circle, fill=blue!85!black, inner sep=1.2pt] at (-0.5, 0) {};
    \node[circle, fill=blue!85!black, inner sep=1.2pt] at (0, 0.5) {};
    \node[circle, fill=blue!85!black, inner sep=1.2pt] at (0, -0.5) {};

    \node[below right=1pt, font=\footnotesize] at (0.5, 0) {$\frac{1}{2}$};
    \node[below left=1pt, font=\footnotesize] at (-0.5, 0) {$-\frac{1}{2}$};
    \node[above left=1pt, font=\footnotesize] at (0, 0.5) {$\frac{1}{2}$};
    \node[below left=1pt, font=\footnotesize] at (0, -0.5) {$-\frac{1}{2}$};

    % 境界方程式の注釈
    \node[right=2pt, font=\footnotesize, text=blue!85!black] at (0.4, 0.35)
    {$(1-a)(1-b) = \frac{1}{2}$};
  \end{tikzpicture}
  \caption{点 $P(a, b)$ の存在範囲（斜線部，境界を含む）}
  \label{fig:p_region}
\end{figure}

従って求める面積$U$は，対称性から
\begin{align}
  U&=4\int_0^{1/2}\left(1-\dfrac{1}{2(1-a)}\right)\,da \\
  &=4\left[a+\dfrac{1}{2}\log |1-a|\right]_0^{1/2} \\
  &=2(1-\log 2)
\end{align}
である．$\cdots$(答)

\newpage
\end{document}
